\newpage
\section{INTRODUCCIÓN}

Este proyecto se centra en el desarrollo de un dispensador de comida y agua para mascotas con capacidades IoT, el proyecto ha sido un esfuerzo integral para crear una solución automatizada y eficiente para el cuidado de las mascotas. Este sistema utiliza sensores ultrasónicos para monitorear los niveles de comida y agua, y permite a los propietarios controlar y supervisar el dispensador a través de una aplicación móvil, utilizando la plataforma Blynk.

Durante el desarrollo del proyecto, se realizaron ajustes significativos en el diseño y la selección de componentes para cumplir con las prioridades definidas. La adaptación a las limitaciones técnicas, como el cambio de motores y la búsqueda de alternativas, permitió avanzar en el desarrollo del dispensador, logrando así el objetivo principal del proyecto.

A lo largo del proceso, se enfrentaron varios problemas técnicos, como la falta de potencia del ESP8266 y dificultades con Arduino Cloud. Para superar estos obstáculos, se implementaron estrategias efectivas, incluyendo la transición a la plataforma Blynk, lo que garantizó el funcionamiento del sistema de dispensación de comida y agua. Esto demostró una notable flexibilidad y capacidad de resolución de problemas.

Se logró implementar funcionalidades básicas de IoT mediante la integración de Blynk, permitiendo el monitoreo remoto del estado de los recipientes y el control del sistema desde un dispositivo móvil. Esta implementación cumplió parcialmente con el objetivo de automatización del suministro, facilitando el control a distancia del sistema.

La construcción del hardware, que incluye el dispensador de comida, el sistema de bombas para agua y los sensores ultrasónicos, fue completada exitosamente. Este aspecto del proyecto proporcionó una base sólida para el funcionamiento básico del sistema, permitiendo que el dispensador cumpliera con sus funciones principales de manera efectiva.

Para consultar el trabajo realizado, puede consultarse el siguiente repositorio de trabajo, en la rama llamada ``Proj'': \url{https://github.com/SaboEIE/IE-0624_IS_2024_B87103/tree/Proj}. El código principal en C (archivo .ino), puede consultarse en la ruta \path{Proj/src}. 

\section{JUSTIFICACIÓN}
Este proyecto busca mejorar el bienestar animal al garantizar que las mascotas y animales de granja reciban los nutrientes necesarios en los momentos adecuados. Además, ofrece una solución conveniente para los propietarios, especialmente aquellos con horarios ocupados, al automatizar el proceso de alimentación y liberarlos de la preocupación por mantener horarios específicos de alimentación. El uso de tecnología conectada a Internet permite el monitoreo remoto de la alimentación de los animales, lo que resulta especialmente útil y permite a los propietarios estar ausentes durante largos períodos de tiempo. Además, la automatización del suministro de comida y agua facilita el control preciso del consumo, evitando el desperdicio y garantizando un suministro adecuado sin excesos. 

\section{OBJETIVOS}

\subsection{Objetivo general}
Diseñar y desarrollar un dispositivo suplidor inteligente de agua y comida para animales utilizando ESP8266, con el fin de mejorar el bienestar animal y proporcionar una solución conveniente para los propietarios, automatizando el proceso de alimentación y permitiendo el monitoreo remoto de la alimentación de los animales.

\subsection{Objetivos específicos}

\begin{enumerate}
    \item Implementar un sistema de suministro de agua y comida que permita tanto la solicitud directa del usuario como una rutina automatizada con tiempos predefinidos, utilizando tecnología ESP8266 para controlar los mecanismos de dispensación.
    \item Integrar la aplicación Arduino Cloud para enviar notificaciones al teléfono móvil del usuario, informando sobre el contenido de los recipientes y los tanques principales de agua y comida, así como para permitir la programación de suministros automáticos y solicitudes adrede.
    \item Garantizar un control preciso del consumo de comida y agua mediante la automatización del suministro, evitando el desperdicio y asegurando un suministro adecuado sin excesos. 
\end{enumerate}
 